\documentclass[11pt,a4paper]{article}

\usepackage[utf8]{inputenc}
\usepackage[T1]{fontenc}
\usepackage[french]{babel}
\usepackage[margin=2.5cm]{geometry}
\usepackage{enumitem}
\usepackage{booktabs}
\usepackage{array}
\usepackage{xcolor}
\usepackage{titlesec}
\usepackage{fancyhdr}

\definecolor{accent}{HTML}{CC02A4}
\definecolor{soft}{HTML}{555555}

\titleformat{\section}{\Large\bfseries\color{accent}}{\thesection}{1em}{}
\titleformat{\subsection}{\large\bfseries}{\thesubsection}{1em}{}

\pagestyle{fancy}
\fancyhf{}
\rhead{\textcolor{soft}{GoPro Gimbal Bridge}}
\lhead{\textcolor{soft}{Mode d'emploi}}
\cfoot{\thepage}

\newcommand{\btn}[1]{\textbf{[#1]}}

\title{\textbf{\textcolor{accent}{GoPro Gimbal Bridge}}\\[0.3em]
       \large Mode d'emploi}
\author{}
\date{}

\begin{document}

\maketitle
\thispagestyle{fancy}

\noindent
Ce système permet de piloter un nacelle (gimbal) motorisée à deux axes et une
caméra GoPro depuis une interface web, ainsi que via des boutons physiques.
Il fournit le contrôle des mouvements, le déclenchement photo/vidéo, la
visualisation du flux vidéo en direct et l'accès à la galerie de la caméra.

\tableofcontents
\vspace{1em}

% =====================================================================
\section{Mise en route}

\begin{enumerate}[leftmargin=*]
  \item \textbf{Alimenter le système.} À l'allumage, les LEDs affichent une
        animation arc-en-ciel, puis l'état du système.
  \item \textbf{Se connecter au réseau WiFi} émis par le boîtier
        (par défaut : \texttt{GoPro-Bridge}). Le mot de passe est défini dans
        l'onglet Configuration.
  \item \textbf{Ouvrir la page} dans un navigateur :
        \texttt{https://192.168.4.1}
  \item \textbf{Premier appareil :} un avertissement « Connexion non sécurisée »
        apparaît. Cliquer sur \emph{Avancé} puis \emph{Continuer}. Pour le faire
        disparaître définitivement, télécharger et installer le certificat
        (voir section Configuration).
\end{enumerate}

\noindent
La connexion à la caméra (Bluetooth puis WiFi) se fait automatiquement au
démarrage. Comptez environ 30 secondes avant que tout soit opérationnel.

% =====================================================================
\section{Les boutons physiques}

Quatre boutons commandent les moteurs, un bouton commande le déclencheur.

\begin{center}
\begin{tabular}{>{\raggedright}p{4.5cm} p{9.5cm}}
\toprule
\textbf{Action} & \textbf{Effet} \\
\midrule
\btn{M1 +} / \btn{M1 --} & Déplace le moteur de l'axe 1 dans un sens ou l'autre. \\
\btn{M2 +} / \btn{M2 --} & Déplace le moteur de l'axe 2 dans un sens ou l'autre. \\
\btn{Déclencheur} & Prend une photo, ou démarre / arrête un enregistrement vidéo selon le mode actif. \\
\addlinespace
\btn{M1 +} \emph{et} \btn{M1 --} ensemble & Réactive le moteur 1 s'il a été coupé (sécurité). Il se recentre alors automatiquement. \\
\btn{M2 +} \emph{et} \btn{M2 --} ensemble & Idem pour le moteur 2. \\
\addlinespace
\textbf{Les 4 boutons moteur maintenus 3 s} & Lance la \emph{calibration} des butées (voir section dédiée). \\
\bottomrule
\end{tabular}
\end{center}

\noindent
Le comportement des boutons moteur (pas-à-pas ou continu) et la vitesse se
règlent depuis l'interface web (onglet Contrôle).

% =====================================================================
\section{Les voyants lumineux (LEDs)}

La bande comporte 7 LEDs qui renseignent en un coup d'\oe il sur l'état du
système.

\begin{center}
\begin{tabular}{c >{\raggedright}p{4cm} p{8.5cm}}
\toprule
\textbf{LED} & \textbf{Rôle} & \textbf{Signification} \\
\midrule
1 & Bluetooth (caméra) & Vert = connectée \quad Rouge = non connectée \\
2 & Réseau WiFi & Vert = un appareil est connecté au boîtier \\
3 & État moteurs & Éteinte = tout va bien \quad Jaune = moteur 1 coupé \quad Bleu = moteur 2 coupé \quad Rouge = les deux coupés \\
4 & Prise de vue & Flash blanc = photo prise \quad Rouge clignotant = enregistrement vidéo en cours \\
5--7 & Batterie & Du vert (pleine) vers le rouge (faible) \\
\bottomrule
\end{tabular}
\end{center}

\noindent
\textbf{Cas particuliers :}
\begin{itemize}[leftmargin=*]
  \item \textbf{Toutes les LEDs en rose} : une calibration est en cours.
  \item \textbf{Toutes les LEDs rouges clignotantes} : batterie critique, recharger
        ou remplacer rapidement.
\end{itemize}

% =====================================================================
\section{L'interface web}

L'interface se compose de quatre onglets accessibles en haut de l'écran :
\textbf{Contrôle}, \textbf{Galerie}, \textbf{Live} et \textbf{Config}.

% ---------------------------------------------------------------------
\subsection{Onglet Contrôle}

\begin{itemize}[leftmargin=*]
  \item \textbf{Bouton principal (déclencheur)} : son libellé s'adapte au mode.
        \begin{itemize}
          \item En mode Photo : \emph{PHOTO} — prend une photo à chaque appui.
          \item En mode Vidéo : \emph{RECORD} pour démarrer, \emph{STOP} pour
                arrêter (le bouton devient vert et clignote pendant
                l'enregistrement).
        \end{itemize}
  \item \textbf{Mode preset} : choix entre \emph{Vidéo} et \emph{Photo}. Le mode
        sélectionné est conservé et réappliqué automatiquement à la caméra.
  \item \textbf{Boutons moteur} : choix du comportement des boutons physiques.
        \begin{itemize}
          \item \emph{Pas-à-pas} : chaque appui déplace le moteur d'un pas fixe.
          \item \emph{Continu} : le moteur bouge tant que le bouton est maintenu.
        \end{itemize}
  \item \textbf{Curseur « Pas »} : règle l'amplitude d'un pas (en degrés) en mode
        pas-à-pas.
  \item \textbf{Curseur « Vitesse »} : règle la vitesse (en degrés par seconde) en
        mode continu.
  \item \textbf{Sleep} : met la caméra en veille.
\end{itemize}

% ---------------------------------------------------------------------
\subsection{Onglet Galerie}

Affiche les photos et vidéos présentes sur la caméra.

\begin{itemize}[leftmargin=*]
  \item \textbf{Actualiser} : recharge la liste des médias.
  \item \textbf{Vignettes} : les aperçus se chargent progressivement. Cliquer sur
        une photo pour l'ouvrir en grand.
  \item \textbf{Dans l'aperçu} : \emph{Télécharger} pour récupérer la photo, ou
        \emph{Supprimer} pour l'effacer de la caméra.
  \item \textbf{Supprimer} (barre d'outils) : efface les éléments sélectionnés.
\end{itemize}

\noindent
\textit{Remarque : le chargement des vignettes peut prendre quelques secondes,
c'est normal. La galerie nécessite que le voyant WiFi soit vert.}

% ---------------------------------------------------------------------
\subsection{Onglet Live}

Affiche le flux vidéo en direct de la caméra via la carte de capture HDMI.

\begin{itemize}[leftmargin=*]
  \item \textbf{Source vidéo} : menu déroulant pour choisir l'entrée vidéo. La
        carte de capture est sélectionnée automatiquement si elle est détectée.
  \item \textbf{Démarrer} / \textbf{Arrêter} : lance ou stoppe l'affichage du flux.
  \item \textbf{Plein écran} (icône en haut à droite de la vidéo) : passe en
        plein écran. Sur tablette/téléphone, toucher l'écran fait apparaître le
        bouton de sortie ; on peut aussi balayer vers le bas.
\end{itemize}

\noindent
\textit{Brancher la carte de capture HDMI avant d'ouvrir cet onglet. Le passage
en plein écran réaffirme également le mode (photo ou vidéo) auprès de la caméra.}

% ---------------------------------------------------------------------
\subsection{Onglet Config}

Permet de paramétrer les connexions et les préférences.

\begin{itemize}[leftmargin=*]
  \item \textbf{WiFi GoPro} : nom (SSID), mot de passe et adresse IP de la caméra.
  \item \textbf{Hotspot Bridge} : nom et mot de passe du réseau WiFi émis par le
        boîtier.
  \item \textbf{Enregistrer} : sauvegarde la configuration. \emph{Un redémarrage
        est nécessaire pour appliquer les changements réseau.}
  \item \textbf{Source vidéo par défaut} : choisit la caméra / carte de capture
        utilisée automatiquement à l'ouverture de l'onglet Live.
  \item \textbf{Télécharger le certificat} : pour installer le certificat de
        sécurité sur un nouvel appareil et supprimer l'avertissement du
        navigateur.
\end{itemize}

% =====================================================================
\section{Calibration des butées}

La calibration détermine les limites de débattement de chaque axe pour éviter de
forcer contre les butées mécaniques.

\begin{enumerate}[leftmargin=*]
  \item \textbf{Maintenir les 4 boutons moteur} (M1+, M1--, M2+, M2--)
        simultanément pendant 3 secondes.
  \item Toutes les LEDs passent au \textbf{rose} : la calibration est en cours.
        Les moteurs se déplacent lentement jusqu'à détecter chaque butée.
  \item À la fin, les limites sont \textbf{enregistrées} et conservées même après
        extinction. Les LEDs reviennent à l'affichage normal.
\end{enumerate}

\noindent
\textit{Des valeurs par défaut sont utilisées tant qu'aucune calibration n'a été
effectuée. Il n'est généralement pas nécessaire de recalibrer sauf en cas de
modification mécanique.}

% =====================================================================
\section{Sécurités automatiques}

\begin{itemize}[leftmargin=*]
  \item \textbf{Protection surcourant} : si un moteur force trop longtemps
        (obstacle, butée), il est automatiquement coupé pour le protéger. La LED
        d'état moteur l'indique (jaune ou bleu).
  \item \textbf{Réactivation} : appuyer simultanément sur les deux boutons d'un
        même axe (ou couper/rallumer depuis le système). Le moteur se
        \textbf{recentre} alors au milieu de sa plage pour se dégager.
  \item \textbf{Limites logicielles} : les moteurs ne dépassent jamais les butées
        calibrées.
\end{itemize}

% =====================================================================
\section{Dépannage rapide}

\begin{center}
\begin{tabular}{>{\raggedright}p{6cm} p{8cm}}
\toprule
\textbf{Symptôme} & \textbf{Solution} \\
\midrule
Le voyant Bluetooth reste rouge & Vérifier que la caméra est allumée et à portée. La connexion se rétablit automatiquement. \\
La galerie n'affiche rien & Attendre que le voyant WiFi soit vert (\textasciitilde30 s après l'allumage). \\
Un moteur ne bouge plus & Il a probablement été coupé par sécurité : réactiver avec les deux boutons de l'axe. \\
Le flux vidéo est noir & Vérifier le branchement HDMI, puis Arrêter et Démarrer à nouveau dans l'onglet Live. \\
Avertissement de sécurité du navigateur & Installer le certificat (onglet Config). \\
Le mode affiché ne correspond pas & Recharger la page : elle se resynchronise sur le mode réel de la caméra. \\
\bottomrule
\end{tabular}
\end{center}

\end{document}
